% =====================================================================
% DRAFT section (26 Nov 2026): identifiability at the matched width, why the
% entropy step fails with finite data, and the chart construction.
% Use: % =====================================================================
% DRAFT section (26 Nov 2026): identifiability at the matched width, why the
% entropy step fails with finite data, and the chart construction.
% Use: % =====================================================================
% DRAFT section (26 Nov 2026): identifiability at the matched width, why the
% entropy step fails with finite data, and the chart construction.
% Use: % =====================================================================
% DRAFT section (26 Nov 2026): identifiability at the matched width, why the
% entropy step fails with finite data, and the chart construction.
% Use: \input{theory-matched-width-draft} after the Background section of draft.tex.
% Needs amsmath, amsthm, booktabs, natbib, xcolor (all loaded by draft.tex).
% The source of every number is in a % comment: items and stages of
% out/burden_feasibility_20261102/HANDOFF_archive_20261127.md (the full record,
% both sessions' copies merged); HANDOFF.md beside it summarises them.
% One new bibliography entry is at the end of this file, commented out.
% =====================================================================
\makeatletter
\@ifundefined{theorem}{\newtheorem{theorem}{Theorem}}{}
\@ifundefined{proposition}{\newtheorem{proposition}{Proposition}}{}
\@ifundefined{lemma}{\newtheorem{lemma}{Lemma}}{}
\makeatother
\providecommand{\R}{\mathbb{R}}
\providecommand{\pending}[1]{{\color{red}\textbf{[PENDING: #1]}}}
\providecommand{\todo}[1]{{\color{blue}\small[TODO: #1]}}

\section{Identifiability at the matched width, and how to reach it}
\label{sec:theory}

Multiview identifiability results promise that a contrastive code of the right
width is a one-to-one function of the content the views share. On our benchmark
that promise is kept in half of the training runs. This section locates the gap
and closes it. We state the result as two steps, each certifying one half of its
conclusion (\S\ref{sec:theory-result}); show that the step certifying
completeness is the one that fails with finite data, and why
(\S\ref{sec:theory-gap}); give a construction that reaches the theorem's
conclusion with a certificate that can be checked on held-out data
(\S\ref{sec:theory-chart}); and state what that certificate does not establish
(\S\ref{sec:theory-check}).

\subsection{Setting}
\label{sec:theory-setting}

Each subject has shared content $c \in C \subset \R^{n_c}$ and, for each view
$v \in \{1, 2\}$ (T1 and FLAIR), a view-specific style $s_v$, drawn
independently given $c$. The views are $x^{(v)} = f_v(c, s_v)$ with $f_v$
smooth and injective. On the fixed-position benchmark the content is seven
anatomical factors (brain volume, ventricle volume, lesion burden, cortical
thickness, temporal volume, left--right asymmetry and sulcal width), so
$n_c = 7$; the style is each view's intensity gain, bias and noise level. Each
view has its own encoder $g_v$, and the two views' codes are compared by the
cosine, that is, on the unit sphere $S^{d-1}$ of a $d$-unit code. With many
negatives, InfoNCE at temperature $\tau$ approaches the alignment
$\mathbb{E}\,\lVert g_1(x^{(1)}) - g_2(x^{(2)}) \rVert^2$ minus the entropy of
each view's code on the sphere \citep{wang2020alignment,zimmermann2021contrastive}.

\subsection{The multiview result is two certificates}
\label{sec:theory-result}

\begin{theorem}[informal; \citealp{vonkugelgen2021self,daunhawer2023identifiability,yao2024multiview}]
\label{thm:block}
Let $g_1$ and $g_2$ minimise the alignment minus the entropy, with the code's
sphere of the content's dimension, $d - 1 = n_c$. Then each view's code is
$g_v(x^{(v)}) = h(c)$ almost surely, for one smooth bijection $h$ of $C$ onto its
image: the code is a function of the content alone, and the content can be
recovered from it.
\end{theorem}
\todo{state the exact theorem and its assumptions as in \citet{yao2024multiview};
here both views observe all of the content, so their partial observability
reduces to the two-view case.}

The proof has two steps, and only the second uses the width.
\begin{enumerate}
\item \emph{Alignment certifies ``only''.} At zero alignment loss
  $g_1(f_1(c, s_1)) = g_2(f_2(c, s_2))$ almost surely. Because $s_1$ and $s_2$
  vary independently given $c$ (with full support, and the maps smooth), neither
  side can depend on its style, so both are one function $h(c)$. This holds at
  any width.
\item \emph{Maximum entropy at the matched width certifies ``all''.} On
  $S^{n_c}$ the entropy is maximal only for the uniform density, which is
  regular. A lemma of \citet{zimmermann2021contrastive} then makes $h$
  one-to-one: a smooth map between simply connected manifolds of the same
  dimension that carries a regular density to a regular density is a
  bijection. \todo{check the lemma's exact statement and number}
\end{enumerate}
Under the cosine the matched width is therefore $d = n_c + 1$ units, eight here;
with an output in $\R^{d}$ rather than on a sphere it would be $n_c$. Above the
matched width, step 2 has nothing to say: an $n_c$-dimensional image in a
higher-dimensional sphere has no finite entropy, so for a wide code the theorem
guarantees only step 1's conclusion.

\subsection{The entropy step fails at the matched width with finite data}
\label{sec:theory-gap}

The entropy in step 2 is never observed. InfoNCE estimates it from batches of
subjects at temperature $\tau$, so it penalises codes that crowd subjects
together at roughly that scale. Consider a code that ignores one direction of
$C$, the variation of one factor's own part, and bends a curved function of the
factors it does hold through the dimension this frees. In the population the
lemma excludes it: its image has dimension $n_c - 1$, a set of measure zero in
$S^{n_c}$, so its density is singular and its entropy is $-\infty$. At the
resolution a batch can see, the bent image can be spread as evenly as a code
that keeps the direction.

Alignment then prefers the curve. A curved function of well-read factors is
computed identically from both views, while the weakest own part (here in the
plane of ventricle size and lesion burden, the variation their causes do not
explain) is determined only partially by each view, so keeping it costs
alignment. The finite-sample optimum can therefore trade the weakest direction
for a curve, and nothing in the objective tells the two apart.

The benchmark shows this trade. % HANDOFF items 46, 56, 58, 61, 63, 64
\begin{itemize}
\item At the matched width (eight units, 2{,}000 training subjects), the code is
  complete\footnote{A code is \emph{complete} when every direction of the seven
  factors' own parts (each factor's variation not explained by its causes) is
  recovered from it with $R^2 \geq .40$ in both views, the ventricle--burden
  plane included. \todo{align with the evaluation section's definition}} in 5
  of 10 training runs that differ only in their seed. % item 61: A100s, seeds 6-15
  Complete codes spend their eighth direction on a curve; incomplete ones spend
  two directions on curves and lose an own-part direction, typically read at
  $R^2$ near zero. % items 46, 56, 58
\item No code carries style at any width: the gain, bias and noise level of
  either view read $R^2 = 0.00$ from codes of 4 to 16 units, linearly and by
  gradient boosting. Step 1 works. % item 63
\item Wider codes put their extra directions into curves. Past the seventh,
  every direction is a joint non-linear function of the seven factors (boosted
  $R^2 \approx .75$) rather than new content, and the codes' intrinsic dimension
  stays at the factors': $6.36$ for 16-unit codes against $6.37$ for the true
  factors on the same subjects (TwoNN). % items 63, 64
\item Nine changes to training at eight units fail to raise the rate: the
  interval at which code batches are rebuilt, the prediction weight, a
  compressed start, a harmonic penalty, the temperature, a projector head at the
  code's width, a wider projector that lets the loss bend the code into sixteen
  outputs (2 of 10 complete), a term penalising what the views still share
  beyond the code (0 of 10), and refitting the head on a frozen backbone.
  % stages 60, 62-64, 70, 71, 71a (item 58); stage 83 (item 69); stage 77 (item 67); stage 80 (item 65)
\item Training longer helps slowly and not always. With 40{,}000 further steps,
  three times the usual length, one of two incomplete codes recovered its lost
  direction (its $R^2$ rose steadily from $.01$ to $.44$), while the other's
  stayed at zero; 10{,}000 further steps had changed neither. The weakest factor
  is learned last.
  % stage 81 (item 70); stage 65 for the 10,000-step continuation
\end{itemize}
The content is identifiable at the matched width, but with finite data the
objective that should find it does not certify completeness.

\subsection{Identify wide, chart narrow}
\label{sec:theory-chart}

The two certificates can be separated. Step 1 does not need the matched width,
and completeness can be certified by something other than entropy. We
therefore (i) train the same objective at a width slightly above the matched
one, where the curves have room and step 1 still holds, and (ii) bring that code
down to the content's dimension with a non-linear chart fitted to rebuild it.

\paragraph{Construction (no labels).}
\begin{enumerate}
\item Train the encoders at a width $d_w > n_c + 1$ (ten units here). By step 1,
  $g_v(x^{(v)}) = \varphi(c)$ for a content-only $\varphi$; that $\varphi$ is
  one-to-one is measured, not guaranteed.
\item In each view, standardise the code, fit a canonical correlation analysis
  between the views on training subjects, keep the directions whose agreement
  on held-out training subjects exceeds $.5$ and a shuffled-pairing null, and
  whiten them. Write $u^{(v)} \in \R^{k}$ for the result.
\item Fit an encoder $\psi : \R^{k} \to \R^{m}$ and a decoder
  $\psi' : \R^{m} \to \R^{k}$ (small MLPs, shared by the views) to minimise
  $\mathbb{E}\,\lVert \psi'(\psi(u^{(v)})) - u^{(v)} \rVert^2$, with early
  stopping on held-out training subjects and the best of five restarts. The
  narrow code is $z^{(v)} = \psi(u^{(v)})$.
\end{enumerate}

\begin{proposition}[chart]
\label{prop:chart}
Let $\varphi : C \to \R^{k}$ be continuous and one-to-one, and let
$\psi'(\psi(u)) = u$ for every $u \in \varphi(C)$. Then $\psi \circ \varphi$ is
continuous and one-to-one. If $m = n_c$, it is a homeomorphism of $C$ onto an
open subset of $\R^{n_c}$, and a diffeomorphism when $\psi$ and $\varphi$ are
smooth with Jacobians of full rank.
\end{proposition}
\begin{proof}
If $\psi(\varphi(c)) = \psi(\varphi(c'))$, applying $\psi'$ gives
$\varphi(c) = \varphi(c')$, so $c = c'$. Continuity is inherited. For
$m = n_c$, a continuous one-to-one map from an open subset of $\R^{n_c}$ into
$\R^{n_c}$ is open by invariance of domain, hence a homeomorphism onto its
image.
\end{proof}

The narrow code is the same function of either view's wide code, so it inherits
step 1's conclusion, and Proposition~\ref{prop:chart} supplies step 2's.
Together they give the conclusion of Theorem~\ref{thm:block} at width $n_c$, one
unit below the cosine's matched width, because the chart's output is not on a
sphere.

Reconstruction succeeds where entropy did not for two reasons. The chart is
fitted to content that is already aligned, so keeping the weakest direction
costs it nothing. And with whitened targets every shared direction counts
equally, while a curve needs no unit of its own: the decoder rebuilds it from
the coordinates of the factors it is a function of. Entropy has neither
property.

\begin{table}[t]
\centering
\caption{Charts of wide codes on the fixed-position benchmark. A code counts when
it is complete. The linear chart (the top principal components of the same
whitened variates) is the best linear compression of the same width.}
\label{tab:chart}
\begin{tabular}{lccc}
\toprule
Wide code $\to$ chart & wide codes & non-linear chart & linear chart\\
\midrule
10 $\to$ 8 units & 4 of 5 & 4 of 4 & 0 of 4\\
10 $\to$ 7 units & 4 of 5 & 4 of 4 & 0 of 4\\
9 $\to$ 8 (code batches every 500 steps) & 3 of 4 & 3 of 3 & 0 of 3\\
9 $\to$ 8 (code batches every 100 steps) & 1 of 4 & 1 of 1 & 0 of 1\\
12 $\to$ 8 & 3 of 3 & 2 of 3 & 0 of 3\\
16 $\to$ 8 & 2 of 3 & 0 of 2 & 0 of 2\\
\midrule
Trained at 8 units (reference) & \multicolumn{3}{c}{5 of 10 complete}\\
\bottomrule
\end{tabular}
\end{table}
% HANDOFF item 66 (stages 82 and 82b); item 61 for the reference row

Wherever its wide code is complete, the chart of a 9- or 10-unit code is
complete: 8 of 8 at eight units and 4 of 4 at seven (Table~\ref{tab:chart}). Its
own-part $R^2$ stays within $.05$ of the wide code's (for example ventricle
$.80 \to .78$ and sulcal width $.74 \to .73$, means over the complete 10-unit
codes and both views). The linear chart is never complete. The charts also keep
their wide codes' causal outputs: average treatment effects estimated with a
smooth learner agree within $.02$, and causal discovery with the code as one
block node recovers its edges to tabular variables with $F_1$ of $.88$--$.94$,
against $.88$--$.96$ for the wide codes and $.98$ for the true own parts.
% item 66, downstream: causal/stage82/tabular_rbf.txt, block_tabular.txt

\subsection{The rebuild check, and what it does not certify}
\label{sec:theory-check}

Proposition~\ref{prop:chart} needs exact rebuilding on $\varphi(C)$. With finite
data, the rebuild error measures how far from it a chart is.

\begin{lemma}[rebuild bound]
\label{lem:rebuild}
For any $\psi$ and $\psi'$ and every direction $j$ of $u$,
$\mathbb{E}\,\mathrm{Var}\bigl(u_j \mid \psi(u)\bigr) \le
\mathbb{E}\,\bigl(u_j - \psi'(\psi(u))_j\bigr)^2$.
\end{lemma}
\begin{proof}
Among all functions of $\psi(u)$, the conditional mean
$\mathbb{E}[u_j \mid \psi(u)]$ has the smallest mean squared error for $u_j$, and
$\psi'(\psi(u))_j$ is one of those functions.
\end{proof}

The share of each whitened shared direction that the narrow code leaves
undetermined is therefore at most its rebuild error, and exact rebuilding gives
the proposition's hypothesis. We estimate the error per direction on test
subjects that neither the canonical analysis, the chart nor its early stopping
saw, with a 95\% bootstrap upper bound over subjects. For the complete 10-unit
codes, the worst direction of the 8-unit charts is at most $.017$ (the mean over
directions at most $.004$), and that of the 7-unit charts at most $.033$.
% item 66, certificate: causal/stage82/cert_u{7,8}{a,b}.txt

Three properties of the check matter.
\begin{itemize}
\item \emph{It is per whitened direction.} An average error, or one in
  unwhitened units, is dominated by high-variance structure, so a faint factor
  can be missing entirely while the average error looks excellent. This is why
  rebuilding the images would be a weak test of whether a faint factor is held.
\item \emph{It is on held-out subjects.} On the fitting subjects near-perfect
  rebuilding is cheap: almost any map is one-to-one on a finite set, and a
  flexible decoder can memorise it.
\item \emph{It certifies information, not readability, and only for continuous
  charts.} A one-to-one code can hold a factor in a form a probe cannot read,
  and without continuity a single real number can encode several factors, so
  the dimension count in Proposition~\ref{prop:chart} rests on continuity. The
  own-part readouts check readability directly.
\end{itemize}

The check does not certify the wide code. It compares the chart with its wide
code, not the wide code with the content, and a chart cannot add what its wide
code lacks: the one incomplete 10-unit code is rebuilt to within $.036$, and its
chart is still incomplete. The construction is therefore as reliable as the wide
code's completeness, 4 of 5 at ten units. Wider codes chart worse with finite
data (2 of 3 at twelve units, 0 of 2 at sixteen): every extra curve direction
counts fully in the whitened target, so the bottleneck again chooses between a
curve and the weakest factor. A slightly wide code is the right parent. Two
label-free steps complete the construction: choosing a complete code among
several wide ones, and setting the chart's width when $n_c$ is unknown, as on
ADNI.

\paragraph{Choosing a complete wide code without labels.} Every content-only
code is a function of the content, so a complete code determines every other
run's code, while a code that misses a part cannot determine one that holds it.
We therefore train several seeds and score each code by how much of the others'
shared content it predicts. For a candidate $A$ and a reference $B$, each view
of $A$ predicts the other view's shared variates of $B$ (so view-specific
variation cannot count), against $B$'s own view as the ceiling; $q(A \mid B)$ is
the smallest share, over the directions of $B$'s whitened variates, of what $B$'s
own view explains that $A$ also explains. A code's score is its minimum $q$ over
the other seeds, and we keep the code with the highest score. At nine and ten
units the complete code scores higher in 10 of 10 complete--incomplete pairs,
where the backbone-leftover criterion of \S\todo{ref} gets 7, and among any
three seeds the top score is complete in every triple that contains a complete
code. At eight units 68 of 69 pairs are right. Because the score reads a code
against other runs' codes rather than against its own backbone, a backbone that
carries less cannot lower it: the 20 codes of a training variant that degraded
its backbone, all incomplete, score below every complete reference code.
% HANDOFF item 68 (stage 84 A); the variant is stage 77 (item 67); rho1 at 9-10 units: item 60

\paragraph{Setting the chart's width without labels.} The rebuild check also
gives the width. Charting each complete 10-unit code at decreasing widths, the
worst-direction error (95\% bound) stays at or below $.033$ down to seven units
and jumps to $.17$--$.31$ at six, where every chart loses an own part (sulcal
width first). The smallest width at which the chart rebuilds the wide code to
within $.10$ is therefore $n_c = 7$ for every complete code, with no labels and
no calibration borrowed from known latents. The one incomplete 10-unit code
shows no jump at seven and is rebuilt at six units to within $.08$: its shared
part is one dimension short, as its own parts confirm, so a width one below the
others is a second sign that a wide code is incomplete.
% HANDOFF item 68 (stage 84 B): causal/stage84/width.txt

% --- add to the bibliography of draft.tex ---------------------------------
% \bibitem[Wang and Isola(2020)]{wang2020alignment}
% T.~Wang and P.~Isola.
% \newblock Understanding contrastive representation learning through alignment
% and uniformity on the hypersphere.
% \newblock \emph{ICML}, 2020.
 after the Background section of draft.tex.
% Needs amsmath, amsthm, booktabs, natbib, xcolor (all loaded by draft.tex).
% The source of every number is in a % comment: items and stages of
% out/burden_feasibility_20261102/HANDOFF_archive_20261127.md (the full record,
% both sessions' copies merged); HANDOFF.md beside it summarises them.
% One new bibliography entry is at the end of this file, commented out.
% =====================================================================
\makeatletter
\@ifundefined{theorem}{\newtheorem{theorem}{Theorem}}{}
\@ifundefined{proposition}{\newtheorem{proposition}{Proposition}}{}
\@ifundefined{lemma}{\newtheorem{lemma}{Lemma}}{}
\makeatother
\providecommand{\R}{\mathbb{R}}
\providecommand{\pending}[1]{{\color{red}\textbf{[PENDING: #1]}}}
\providecommand{\todo}[1]{{\color{blue}\small[TODO: #1]}}

\section{Identifiability at the matched width, and how to reach it}
\label{sec:theory}

Multiview identifiability results promise that a contrastive code of the right
width is a one-to-one function of the content the views share. On our benchmark
that promise is kept in half of the training runs. This section locates the gap
and closes it. We state the result as two steps, each certifying one half of its
conclusion (\S\ref{sec:theory-result}); show that the step certifying
completeness is the one that fails with finite data, and why
(\S\ref{sec:theory-gap}); give a construction that reaches the theorem's
conclusion with a certificate that can be checked on held-out data
(\S\ref{sec:theory-chart}); and state what that certificate does not establish
(\S\ref{sec:theory-check}).

\subsection{Setting}
\label{sec:theory-setting}

Each subject has shared content $c \in C \subset \R^{n_c}$ and, for each view
$v \in \{1, 2\}$ (T1 and FLAIR), a view-specific style $s_v$, drawn
independently given $c$. The views are $x^{(v)} = f_v(c, s_v)$ with $f_v$
smooth and injective. On the fixed-position benchmark the content is seven
anatomical factors (brain volume, ventricle volume, lesion burden, cortical
thickness, temporal volume, left--right asymmetry and sulcal width), so
$n_c = 7$; the style is each view's intensity gain, bias and noise level. Each
view has its own encoder $g_v$, and the two views' codes are compared by the
cosine, that is, on the unit sphere $S^{d-1}$ of a $d$-unit code. With many
negatives, InfoNCE at temperature $\tau$ approaches the alignment
$\mathbb{E}\,\lVert g_1(x^{(1)}) - g_2(x^{(2)}) \rVert^2$ minus the entropy of
each view's code on the sphere \citep{wang2020alignment,zimmermann2021contrastive}.

\subsection{The multiview result is two certificates}
\label{sec:theory-result}

\begin{theorem}[informal; \citealp{vonkugelgen2021self,daunhawer2023identifiability,yao2024multiview}]
\label{thm:block}
Let $g_1$ and $g_2$ minimise the alignment minus the entropy, with the code's
sphere of the content's dimension, $d - 1 = n_c$. Then each view's code is
$g_v(x^{(v)}) = h(c)$ almost surely, for one smooth bijection $h$ of $C$ onto its
image: the code is a function of the content alone, and the content can be
recovered from it.
\end{theorem}
\todo{state the exact theorem and its assumptions as in \citet{yao2024multiview};
here both views observe all of the content, so their partial observability
reduces to the two-view case.}

The proof has two steps, and only the second uses the width.
\begin{enumerate}
\item \emph{Alignment certifies ``only''.} At zero alignment loss
  $g_1(f_1(c, s_1)) = g_2(f_2(c, s_2))$ almost surely. Because $s_1$ and $s_2$
  vary independently given $c$ (with full support, and the maps smooth), neither
  side can depend on its style, so both are one function $h(c)$. This holds at
  any width.
\item \emph{Maximum entropy at the matched width certifies ``all''.} On
  $S^{n_c}$ the entropy is maximal only for the uniform density, which is
  regular. A lemma of \citet{zimmermann2021contrastive} then makes $h$
  one-to-one: a smooth map between simply connected manifolds of the same
  dimension that carries a regular density to a regular density is a
  bijection. \todo{check the lemma's exact statement and number}
\end{enumerate}
Under the cosine the matched width is therefore $d = n_c + 1$ units, eight here;
with an output in $\R^{d}$ rather than on a sphere it would be $n_c$. Above the
matched width, step 2 has nothing to say: an $n_c$-dimensional image in a
higher-dimensional sphere has no finite entropy, so for a wide code the theorem
guarantees only step 1's conclusion.

\subsection{The entropy step fails at the matched width with finite data}
\label{sec:theory-gap}

The entropy in step 2 is never observed. InfoNCE estimates it from batches of
subjects at temperature $\tau$, so it penalises codes that crowd subjects
together at roughly that scale. Consider a code that ignores one direction of
$C$, the variation of one factor's own part, and bends a curved function of the
factors it does hold through the dimension this frees. In the population the
lemma excludes it: its image has dimension $n_c - 1$, a set of measure zero in
$S^{n_c}$, so its density is singular and its entropy is $-\infty$. At the
resolution a batch can see, the bent image can be spread as evenly as a code
that keeps the direction.

Alignment then prefers the curve. A curved function of well-read factors is
computed identically from both views, while the weakest own part (here in the
plane of ventricle size and lesion burden, the variation their causes do not
explain) is determined only partially by each view, so keeping it costs
alignment. The finite-sample optimum can therefore trade the weakest direction
for a curve, and nothing in the objective tells the two apart.

The benchmark shows this trade. % HANDOFF items 46, 56, 58, 61, 63, 64
\begin{itemize}
\item At the matched width (eight units, 2{,}000 training subjects), the code is
  complete\footnote{A code is \emph{complete} when every direction of the seven
  factors' own parts (each factor's variation not explained by its causes) is
  recovered from it with $R^2 \geq .40$ in both views, the ventricle--burden
  plane included. \todo{align with the evaluation section's definition}} in 5
  of 10 training runs that differ only in their seed. % item 61: A100s, seeds 6-15
  Complete codes spend their eighth direction on a curve; incomplete ones spend
  two directions on curves and lose an own-part direction, typically read at
  $R^2$ near zero. % items 46, 56, 58
\item No code carries style at any width: the gain, bias and noise level of
  either view read $R^2 = 0.00$ from codes of 4 to 16 units, linearly and by
  gradient boosting. Step 1 works. % item 63
\item Wider codes put their extra directions into curves. Past the seventh,
  every direction is a joint non-linear function of the seven factors (boosted
  $R^2 \approx .75$) rather than new content, and the codes' intrinsic dimension
  stays at the factors': $6.36$ for 16-unit codes against $6.37$ for the true
  factors on the same subjects (TwoNN). % items 63, 64
\item Nine changes to training at eight units fail to raise the rate: the
  interval at which code batches are rebuilt, the prediction weight, a
  compressed start, a harmonic penalty, the temperature, a projector head at the
  code's width, a wider projector that lets the loss bend the code into sixteen
  outputs (2 of 10 complete), a term penalising what the views still share
  beyond the code (0 of 10), and refitting the head on a frozen backbone.
  % stages 60, 62-64, 70, 71, 71a (item 58); stage 83 (item 69); stage 77 (item 67); stage 80 (item 65)
\item Training longer helps slowly and not always. With 40{,}000 further steps,
  three times the usual length, one of two incomplete codes recovered its lost
  direction (its $R^2$ rose steadily from $.01$ to $.44$), while the other's
  stayed at zero; 10{,}000 further steps had changed neither. The weakest factor
  is learned last.
  % stage 81 (item 70); stage 65 for the 10,000-step continuation
\end{itemize}
The content is identifiable at the matched width, but with finite data the
objective that should find it does not certify completeness.

\subsection{Identify wide, chart narrow}
\label{sec:theory-chart}

The two certificates can be separated. Step 1 does not need the matched width,
and completeness can be certified by something other than entropy. We
therefore (i) train the same objective at a width slightly above the matched
one, where the curves have room and step 1 still holds, and (ii) bring that code
down to the content's dimension with a non-linear chart fitted to rebuild it.

\paragraph{Construction (no labels).}
\begin{enumerate}
\item Train the encoders at a width $d_w > n_c + 1$ (ten units here). By step 1,
  $g_v(x^{(v)}) = \varphi(c)$ for a content-only $\varphi$; that $\varphi$ is
  one-to-one is measured, not guaranteed.
\item In each view, standardise the code, fit a canonical correlation analysis
  between the views on training subjects, keep the directions whose agreement
  on held-out training subjects exceeds $.5$ and a shuffled-pairing null, and
  whiten them. Write $u^{(v)} \in \R^{k}$ for the result.
\item Fit an encoder $\psi : \R^{k} \to \R^{m}$ and a decoder
  $\psi' : \R^{m} \to \R^{k}$ (small MLPs, shared by the views) to minimise
  $\mathbb{E}\,\lVert \psi'(\psi(u^{(v)})) - u^{(v)} \rVert^2$, with early
  stopping on held-out training subjects and the best of five restarts. The
  narrow code is $z^{(v)} = \psi(u^{(v)})$.
\end{enumerate}

\begin{proposition}[chart]
\label{prop:chart}
Let $\varphi : C \to \R^{k}$ be continuous and one-to-one, and let
$\psi'(\psi(u)) = u$ for every $u \in \varphi(C)$. Then $\psi \circ \varphi$ is
continuous and one-to-one. If $m = n_c$, it is a homeomorphism of $C$ onto an
open subset of $\R^{n_c}$, and a diffeomorphism when $\psi$ and $\varphi$ are
smooth with Jacobians of full rank.
\end{proposition}
\begin{proof}
If $\psi(\varphi(c)) = \psi(\varphi(c'))$, applying $\psi'$ gives
$\varphi(c) = \varphi(c')$, so $c = c'$. Continuity is inherited. For
$m = n_c$, a continuous one-to-one map from an open subset of $\R^{n_c}$ into
$\R^{n_c}$ is open by invariance of domain, hence a homeomorphism onto its
image.
\end{proof}

The narrow code is the same function of either view's wide code, so it inherits
step 1's conclusion, and Proposition~\ref{prop:chart} supplies step 2's.
Together they give the conclusion of Theorem~\ref{thm:block} at width $n_c$, one
unit below the cosine's matched width, because the chart's output is not on a
sphere.

Reconstruction succeeds where entropy did not for two reasons. The chart is
fitted to content that is already aligned, so keeping the weakest direction
costs it nothing. And with whitened targets every shared direction counts
equally, while a curve needs no unit of its own: the decoder rebuilds it from
the coordinates of the factors it is a function of. Entropy has neither
property.

\begin{table}[t]
\centering
\caption{Charts of wide codes on the fixed-position benchmark. A code counts when
it is complete. The linear chart (the top principal components of the same
whitened variates) is the best linear compression of the same width.}
\label{tab:chart}
\begin{tabular}{lccc}
\toprule
Wide code $\to$ chart & wide codes & non-linear chart & linear chart\\
\midrule
10 $\to$ 8 units & 4 of 5 & 4 of 4 & 0 of 4\\
10 $\to$ 7 units & 4 of 5 & 4 of 4 & 0 of 4\\
9 $\to$ 8 (code batches every 500 steps) & 3 of 4 & 3 of 3 & 0 of 3\\
9 $\to$ 8 (code batches every 100 steps) & 1 of 4 & 1 of 1 & 0 of 1\\
12 $\to$ 8 & 3 of 3 & 2 of 3 & 0 of 3\\
16 $\to$ 8 & 2 of 3 & 0 of 2 & 0 of 2\\
\midrule
Trained at 8 units (reference) & \multicolumn{3}{c}{5 of 10 complete}\\
\bottomrule
\end{tabular}
\end{table}
% HANDOFF item 66 (stages 82 and 82b); item 61 for the reference row

Wherever its wide code is complete, the chart of a 9- or 10-unit code is
complete: 8 of 8 at eight units and 4 of 4 at seven (Table~\ref{tab:chart}). Its
own-part $R^2$ stays within $.05$ of the wide code's (for example ventricle
$.80 \to .78$ and sulcal width $.74 \to .73$, means over the complete 10-unit
codes and both views). The linear chart is never complete. The charts also keep
their wide codes' causal outputs: average treatment effects estimated with a
smooth learner agree within $.02$, and causal discovery with the code as one
block node recovers its edges to tabular variables with $F_1$ of $.88$--$.94$,
against $.88$--$.96$ for the wide codes and $.98$ for the true own parts.
% item 66, downstream: causal/stage82/tabular_rbf.txt, block_tabular.txt

\subsection{The rebuild check, and what it does not certify}
\label{sec:theory-check}

Proposition~\ref{prop:chart} needs exact rebuilding on $\varphi(C)$. With finite
data, the rebuild error measures how far from it a chart is.

\begin{lemma}[rebuild bound]
\label{lem:rebuild}
For any $\psi$ and $\psi'$ and every direction $j$ of $u$,
$\mathbb{E}\,\mathrm{Var}\bigl(u_j \mid \psi(u)\bigr) \le
\mathbb{E}\,\bigl(u_j - \psi'(\psi(u))_j\bigr)^2$.
\end{lemma}
\begin{proof}
Among all functions of $\psi(u)$, the conditional mean
$\mathbb{E}[u_j \mid \psi(u)]$ has the smallest mean squared error for $u_j$, and
$\psi'(\psi(u))_j$ is one of those functions.
\end{proof}

The share of each whitened shared direction that the narrow code leaves
undetermined is therefore at most its rebuild error, and exact rebuilding gives
the proposition's hypothesis. We estimate the error per direction on test
subjects that neither the canonical analysis, the chart nor its early stopping
saw, with a 95\% bootstrap upper bound over subjects. For the complete 10-unit
codes, the worst direction of the 8-unit charts is at most $.017$ (the mean over
directions at most $.004$), and that of the 7-unit charts at most $.033$.
% item 66, certificate: causal/stage82/cert_u{7,8}{a,b}.txt

Three properties of the check matter.
\begin{itemize}
\item \emph{It is per whitened direction.} An average error, or one in
  unwhitened units, is dominated by high-variance structure, so a faint factor
  can be missing entirely while the average error looks excellent. This is why
  rebuilding the images would be a weak test of whether a faint factor is held.
\item \emph{It is on held-out subjects.} On the fitting subjects near-perfect
  rebuilding is cheap: almost any map is one-to-one on a finite set, and a
  flexible decoder can memorise it.
\item \emph{It certifies information, not readability, and only for continuous
  charts.} A one-to-one code can hold a factor in a form a probe cannot read,
  and without continuity a single real number can encode several factors, so
  the dimension count in Proposition~\ref{prop:chart} rests on continuity. The
  own-part readouts check readability directly.
\end{itemize}

The check does not certify the wide code. It compares the chart with its wide
code, not the wide code with the content, and a chart cannot add what its wide
code lacks: the one incomplete 10-unit code is rebuilt to within $.036$, and its
chart is still incomplete. The construction is therefore as reliable as the wide
code's completeness, 4 of 5 at ten units. Wider codes chart worse with finite
data (2 of 3 at twelve units, 0 of 2 at sixteen): every extra curve direction
counts fully in the whitened target, so the bottleneck again chooses between a
curve and the weakest factor. A slightly wide code is the right parent. Two
label-free steps complete the construction: choosing a complete code among
several wide ones, and setting the chart's width when $n_c$ is unknown, as on
ADNI.

\paragraph{Choosing a complete wide code without labels.} Every content-only
code is a function of the content, so a complete code determines every other
run's code, while a code that misses a part cannot determine one that holds it.
We therefore train several seeds and score each code by how much of the others'
shared content it predicts. For a candidate $A$ and a reference $B$, each view
of $A$ predicts the other view's shared variates of $B$ (so view-specific
variation cannot count), against $B$'s own view as the ceiling; $q(A \mid B)$ is
the smallest share, over the directions of $B$'s whitened variates, of what $B$'s
own view explains that $A$ also explains. A code's score is its minimum $q$ over
the other seeds, and we keep the code with the highest score. At nine and ten
units the complete code scores higher in 10 of 10 complete--incomplete pairs,
where the backbone-leftover criterion of \S\todo{ref} gets 7, and among any
three seeds the top score is complete in every triple that contains a complete
code. At eight units 68 of 69 pairs are right. Because the score reads a code
against other runs' codes rather than against its own backbone, a backbone that
carries less cannot lower it: the 20 codes of a training variant that degraded
its backbone, all incomplete, score below every complete reference code.
% HANDOFF item 68 (stage 84 A); the variant is stage 77 (item 67); rho1 at 9-10 units: item 60

\paragraph{Setting the chart's width without labels.} The rebuild check also
gives the width. Charting each complete 10-unit code at decreasing widths, the
worst-direction error (95\% bound) stays at or below $.033$ down to seven units
and jumps to $.17$--$.31$ at six, where every chart loses an own part (sulcal
width first). The smallest width at which the chart rebuilds the wide code to
within $.10$ is therefore $n_c = 7$ for every complete code, with no labels and
no calibration borrowed from known latents. The one incomplete 10-unit code
shows no jump at seven and is rebuilt at six units to within $.08$: its shared
part is one dimension short, as its own parts confirm, so a width one below the
others is a second sign that a wide code is incomplete.
% HANDOFF item 68 (stage 84 B): causal/stage84/width.txt

% --- add to the bibliography of draft.tex ---------------------------------
% \bibitem[Wang and Isola(2020)]{wang2020alignment}
% T.~Wang and P.~Isola.
% \newblock Understanding contrastive representation learning through alignment
% and uniformity on the hypersphere.
% \newblock \emph{ICML}, 2020.
 after the Background section of draft.tex.
% Needs amsmath, amsthm, booktabs, natbib, xcolor (all loaded by draft.tex).
% The source of every number is in a % comment: items and stages of
% out/burden_feasibility_20261102/HANDOFF_archive_20261127.md (the full record,
% both sessions' copies merged); HANDOFF.md beside it summarises them.
% One new bibliography entry is at the end of this file, commented out.
% =====================================================================
\makeatletter
\@ifundefined{theorem}{\newtheorem{theorem}{Theorem}}{}
\@ifundefined{proposition}{\newtheorem{proposition}{Proposition}}{}
\@ifundefined{lemma}{\newtheorem{lemma}{Lemma}}{}
\makeatother
\providecommand{\R}{\mathbb{R}}
\providecommand{\pending}[1]{{\color{red}\textbf{[PENDING: #1]}}}
\providecommand{\todo}[1]{{\color{blue}\small[TODO: #1]}}

\section{Identifiability at the matched width, and how to reach it}
\label{sec:theory}

Multiview identifiability results promise that a contrastive code of the right
width is a one-to-one function of the content the views share. On our benchmark
that promise is kept in half of the training runs. This section locates the gap
and closes it. We state the result as two steps, each certifying one half of its
conclusion (\S\ref{sec:theory-result}); show that the step certifying
completeness is the one that fails with finite data, and why
(\S\ref{sec:theory-gap}); give a construction that reaches the theorem's
conclusion with a certificate that can be checked on held-out data
(\S\ref{sec:theory-chart}); and state what that certificate does not establish
(\S\ref{sec:theory-check}).

\subsection{Setting}
\label{sec:theory-setting}

Each subject has shared content $c \in C \subset \R^{n_c}$ and, for each view
$v \in \{1, 2\}$ (T1 and FLAIR), a view-specific style $s_v$, drawn
independently given $c$. The views are $x^{(v)} = f_v(c, s_v)$ with $f_v$
smooth and injective. On the fixed-position benchmark the content is seven
anatomical factors (brain volume, ventricle volume, lesion burden, cortical
thickness, temporal volume, left--right asymmetry and sulcal width), so
$n_c = 7$; the style is each view's intensity gain, bias and noise level. Each
view has its own encoder $g_v$, and the two views' codes are compared by the
cosine, that is, on the unit sphere $S^{d-1}$ of a $d$-unit code. With many
negatives, InfoNCE at temperature $\tau$ approaches the alignment
$\mathbb{E}\,\lVert g_1(x^{(1)}) - g_2(x^{(2)}) \rVert^2$ minus the entropy of
each view's code on the sphere \citep{wang2020alignment,zimmermann2021contrastive}.

\subsection{The multiview result is two certificates}
\label{sec:theory-result}

\begin{theorem}[informal; \citealp{vonkugelgen2021self,daunhawer2023identifiability,yao2024multiview}]
\label{thm:block}
Let $g_1$ and $g_2$ minimise the alignment minus the entropy, with the code's
sphere of the content's dimension, $d - 1 = n_c$. Then each view's code is
$g_v(x^{(v)}) = h(c)$ almost surely, for one smooth bijection $h$ of $C$ onto its
image: the code is a function of the content alone, and the content can be
recovered from it.
\end{theorem}
\todo{state the exact theorem and its assumptions as in \citet{yao2024multiview};
here both views observe all of the content, so their partial observability
reduces to the two-view case.}

The proof has two steps, and only the second uses the width.
\begin{enumerate}
\item \emph{Alignment certifies ``only''.} At zero alignment loss
  $g_1(f_1(c, s_1)) = g_2(f_2(c, s_2))$ almost surely. Because $s_1$ and $s_2$
  vary independently given $c$ (with full support, and the maps smooth), neither
  side can depend on its style, so both are one function $h(c)$. This holds at
  any width.
\item \emph{Maximum entropy at the matched width certifies ``all''.} On
  $S^{n_c}$ the entropy is maximal only for the uniform density, which is
  regular. A lemma of \citet{zimmermann2021contrastive} then makes $h$
  one-to-one: a smooth map between simply connected manifolds of the same
  dimension that carries a regular density to a regular density is a
  bijection. \todo{check the lemma's exact statement and number}
\end{enumerate}
Under the cosine the matched width is therefore $d = n_c + 1$ units, eight here;
with an output in $\R^{d}$ rather than on a sphere it would be $n_c$. Above the
matched width, step 2 has nothing to say: an $n_c$-dimensional image in a
higher-dimensional sphere has no finite entropy, so for a wide code the theorem
guarantees only step 1's conclusion.

\subsection{The entropy step fails at the matched width with finite data}
\label{sec:theory-gap}

The entropy in step 2 is never observed. InfoNCE estimates it from batches of
subjects at temperature $\tau$, so it penalises codes that crowd subjects
together at roughly that scale. Consider a code that ignores one direction of
$C$, the variation of one factor's own part, and bends a curved function of the
factors it does hold through the dimension this frees. In the population the
lemma excludes it: its image has dimension $n_c - 1$, a set of measure zero in
$S^{n_c}$, so its density is singular and its entropy is $-\infty$. At the
resolution a batch can see, the bent image can be spread as evenly as a code
that keeps the direction.

Alignment then prefers the curve. A curved function of well-read factors is
computed identically from both views, while the weakest own part (here in the
plane of ventricle size and lesion burden, the variation their causes do not
explain) is determined only partially by each view, so keeping it costs
alignment. The finite-sample optimum can therefore trade the weakest direction
for a curve, and nothing in the objective tells the two apart.

The benchmark shows this trade. % HANDOFF items 46, 56, 58, 61, 63, 64
\begin{itemize}
\item At the matched width (eight units, 2{,}000 training subjects), the code is
  complete\footnote{A code is \emph{complete} when every direction of the seven
  factors' own parts (each factor's variation not explained by its causes) is
  recovered from it with $R^2 \geq .40$ in both views, the ventricle--burden
  plane included. \todo{align with the evaluation section's definition}} in 5
  of 10 training runs that differ only in their seed. % item 61: A100s, seeds 6-15
  Complete codes spend their eighth direction on a curve; incomplete ones spend
  two directions on curves and lose an own-part direction, typically read at
  $R^2$ near zero. % items 46, 56, 58
\item No code carries style at any width: the gain, bias and noise level of
  either view read $R^2 = 0.00$ from codes of 4 to 16 units, linearly and by
  gradient boosting. Step 1 works. % item 63
\item Wider codes put their extra directions into curves. Past the seventh,
  every direction is a joint non-linear function of the seven factors (boosted
  $R^2 \approx .75$) rather than new content, and the codes' intrinsic dimension
  stays at the factors': $6.36$ for 16-unit codes against $6.37$ for the true
  factors on the same subjects (TwoNN). % items 63, 64
\item Nine changes to training at eight units fail to raise the rate: the
  interval at which code batches are rebuilt, the prediction weight, a
  compressed start, a harmonic penalty, the temperature, a projector head at the
  code's width, a wider projector that lets the loss bend the code into sixteen
  outputs (2 of 10 complete), a term penalising what the views still share
  beyond the code (0 of 10), and refitting the head on a frozen backbone.
  % stages 60, 62-64, 70, 71, 71a (item 58); stage 83 (item 69); stage 77 (item 67); stage 80 (item 65)
\item Training longer helps slowly and not always. With 40{,}000 further steps,
  three times the usual length, one of two incomplete codes recovered its lost
  direction (its $R^2$ rose steadily from $.01$ to $.44$), while the other's
  stayed at zero; 10{,}000 further steps had changed neither. The weakest factor
  is learned last.
  % stage 81 (item 70); stage 65 for the 10,000-step continuation
\end{itemize}
The content is identifiable at the matched width, but with finite data the
objective that should find it does not certify completeness.

\subsection{Identify wide, chart narrow}
\label{sec:theory-chart}

The two certificates can be separated. Step 1 does not need the matched width,
and completeness can be certified by something other than entropy. We
therefore (i) train the same objective at a width slightly above the matched
one, where the curves have room and step 1 still holds, and (ii) bring that code
down to the content's dimension with a non-linear chart fitted to rebuild it.

\paragraph{Construction (no labels).}
\begin{enumerate}
\item Train the encoders at a width $d_w > n_c + 1$ (ten units here). By step 1,
  $g_v(x^{(v)}) = \varphi(c)$ for a content-only $\varphi$; that $\varphi$ is
  one-to-one is measured, not guaranteed.
\item In each view, standardise the code, fit a canonical correlation analysis
  between the views on training subjects, keep the directions whose agreement
  on held-out training subjects exceeds $.5$ and a shuffled-pairing null, and
  whiten them. Write $u^{(v)} \in \R^{k}$ for the result.
\item Fit an encoder $\psi : \R^{k} \to \R^{m}$ and a decoder
  $\psi' : \R^{m} \to \R^{k}$ (small MLPs, shared by the views) to minimise
  $\mathbb{E}\,\lVert \psi'(\psi(u^{(v)})) - u^{(v)} \rVert^2$, with early
  stopping on held-out training subjects and the best of five restarts. The
  narrow code is $z^{(v)} = \psi(u^{(v)})$.
\end{enumerate}

\begin{proposition}[chart]
\label{prop:chart}
Let $\varphi : C \to \R^{k}$ be continuous and one-to-one, and let
$\psi'(\psi(u)) = u$ for every $u \in \varphi(C)$. Then $\psi \circ \varphi$ is
continuous and one-to-one. If $m = n_c$, it is a homeomorphism of $C$ onto an
open subset of $\R^{n_c}$, and a diffeomorphism when $\psi$ and $\varphi$ are
smooth with Jacobians of full rank.
\end{proposition}
\begin{proof}
If $\psi(\varphi(c)) = \psi(\varphi(c'))$, applying $\psi'$ gives
$\varphi(c) = \varphi(c')$, so $c = c'$. Continuity is inherited. For
$m = n_c$, a continuous one-to-one map from an open subset of $\R^{n_c}$ into
$\R^{n_c}$ is open by invariance of domain, hence a homeomorphism onto its
image.
\end{proof}

The narrow code is the same function of either view's wide code, so it inherits
step 1's conclusion, and Proposition~\ref{prop:chart} supplies step 2's.
Together they give the conclusion of Theorem~\ref{thm:block} at width $n_c$, one
unit below the cosine's matched width, because the chart's output is not on a
sphere.

Reconstruction succeeds where entropy did not for two reasons. The chart is
fitted to content that is already aligned, so keeping the weakest direction
costs it nothing. And with whitened targets every shared direction counts
equally, while a curve needs no unit of its own: the decoder rebuilds it from
the coordinates of the factors it is a function of. Entropy has neither
property.

\begin{table}[t]
\centering
\caption{Charts of wide codes on the fixed-position benchmark. A code counts when
it is complete. The linear chart (the top principal components of the same
whitened variates) is the best linear compression of the same width.}
\label{tab:chart}
\begin{tabular}{lccc}
\toprule
Wide code $\to$ chart & wide codes & non-linear chart & linear chart\\
\midrule
10 $\to$ 8 units & 4 of 5 & 4 of 4 & 0 of 4\\
10 $\to$ 7 units & 4 of 5 & 4 of 4 & 0 of 4\\
9 $\to$ 8 (code batches every 500 steps) & 3 of 4 & 3 of 3 & 0 of 3\\
9 $\to$ 8 (code batches every 100 steps) & 1 of 4 & 1 of 1 & 0 of 1\\
12 $\to$ 8 & 3 of 3 & 2 of 3 & 0 of 3\\
16 $\to$ 8 & 2 of 3 & 0 of 2 & 0 of 2\\
\midrule
Trained at 8 units (reference) & \multicolumn{3}{c}{5 of 10 complete}\\
\bottomrule
\end{tabular}
\end{table}
% HANDOFF item 66 (stages 82 and 82b); item 61 for the reference row

Wherever its wide code is complete, the chart of a 9- or 10-unit code is
complete: 8 of 8 at eight units and 4 of 4 at seven (Table~\ref{tab:chart}). Its
own-part $R^2$ stays within $.05$ of the wide code's (for example ventricle
$.80 \to .78$ and sulcal width $.74 \to .73$, means over the complete 10-unit
codes and both views). The linear chart is never complete. The charts also keep
their wide codes' causal outputs: average treatment effects estimated with a
smooth learner agree within $.02$, and causal discovery with the code as one
block node recovers its edges to tabular variables with $F_1$ of $.88$--$.94$,
against $.88$--$.96$ for the wide codes and $.98$ for the true own parts.
% item 66, downstream: causal/stage82/tabular_rbf.txt, block_tabular.txt

\subsection{The rebuild check, and what it does not certify}
\label{sec:theory-check}

Proposition~\ref{prop:chart} needs exact rebuilding on $\varphi(C)$. With finite
data, the rebuild error measures how far from it a chart is.

\begin{lemma}[rebuild bound]
\label{lem:rebuild}
For any $\psi$ and $\psi'$ and every direction $j$ of $u$,
$\mathbb{E}\,\mathrm{Var}\bigl(u_j \mid \psi(u)\bigr) \le
\mathbb{E}\,\bigl(u_j - \psi'(\psi(u))_j\bigr)^2$.
\end{lemma}
\begin{proof}
Among all functions of $\psi(u)$, the conditional mean
$\mathbb{E}[u_j \mid \psi(u)]$ has the smallest mean squared error for $u_j$, and
$\psi'(\psi(u))_j$ is one of those functions.
\end{proof}

The share of each whitened shared direction that the narrow code leaves
undetermined is therefore at most its rebuild error, and exact rebuilding gives
the proposition's hypothesis. We estimate the error per direction on test
subjects that neither the canonical analysis, the chart nor its early stopping
saw, with a 95\% bootstrap upper bound over subjects. For the complete 10-unit
codes, the worst direction of the 8-unit charts is at most $.017$ (the mean over
directions at most $.004$), and that of the 7-unit charts at most $.033$.
% item 66, certificate: causal/stage82/cert_u{7,8}{a,b}.txt

Three properties of the check matter.
\begin{itemize}
\item \emph{It is per whitened direction.} An average error, or one in
  unwhitened units, is dominated by high-variance structure, so a faint factor
  can be missing entirely while the average error looks excellent. This is why
  rebuilding the images would be a weak test of whether a faint factor is held.
\item \emph{It is on held-out subjects.} On the fitting subjects near-perfect
  rebuilding is cheap: almost any map is one-to-one on a finite set, and a
  flexible decoder can memorise it.
\item \emph{It certifies information, not readability, and only for continuous
  charts.} A one-to-one code can hold a factor in a form a probe cannot read,
  and without continuity a single real number can encode several factors, so
  the dimension count in Proposition~\ref{prop:chart} rests on continuity. The
  own-part readouts check readability directly.
\end{itemize}

The check does not certify the wide code. It compares the chart with its wide
code, not the wide code with the content, and a chart cannot add what its wide
code lacks: the one incomplete 10-unit code is rebuilt to within $.036$, and its
chart is still incomplete. The construction is therefore as reliable as the wide
code's completeness, 4 of 5 at ten units. Wider codes chart worse with finite
data (2 of 3 at twelve units, 0 of 2 at sixteen): every extra curve direction
counts fully in the whitened target, so the bottleneck again chooses between a
curve and the weakest factor. A slightly wide code is the right parent. Two
label-free steps complete the construction: choosing a complete code among
several wide ones, and setting the chart's width when $n_c$ is unknown, as on
ADNI.

\paragraph{Choosing a complete wide code without labels.} Every content-only
code is a function of the content, so a complete code determines every other
run's code, while a code that misses a part cannot determine one that holds it.
We therefore train several seeds and score each code by how much of the others'
shared content it predicts. For a candidate $A$ and a reference $B$, each view
of $A$ predicts the other view's shared variates of $B$ (so view-specific
variation cannot count), against $B$'s own view as the ceiling; $q(A \mid B)$ is
the smallest share, over the directions of $B$'s whitened variates, of what $B$'s
own view explains that $A$ also explains. A code's score is its minimum $q$ over
the other seeds, and we keep the code with the highest score. At nine and ten
units the complete code scores higher in 10 of 10 complete--incomplete pairs,
where the backbone-leftover criterion of \S\todo{ref} gets 7, and among any
three seeds the top score is complete in every triple that contains a complete
code. At eight units 68 of 69 pairs are right. Because the score reads a code
against other runs' codes rather than against its own backbone, a backbone that
carries less cannot lower it: the 20 codes of a training variant that degraded
its backbone, all incomplete, score below every complete reference code.
% HANDOFF item 68 (stage 84 A); the variant is stage 77 (item 67); rho1 at 9-10 units: item 60

\paragraph{Setting the chart's width without labels.} The rebuild check also
gives the width. Charting each complete 10-unit code at decreasing widths, the
worst-direction error (95\% bound) stays at or below $.033$ down to seven units
and jumps to $.17$--$.31$ at six, where every chart loses an own part (sulcal
width first). The smallest width at which the chart rebuilds the wide code to
within $.10$ is therefore $n_c = 7$ for every complete code, with no labels and
no calibration borrowed from known latents. The one incomplete 10-unit code
shows no jump at seven and is rebuilt at six units to within $.08$: its shared
part is one dimension short, as its own parts confirm, so a width one below the
others is a second sign that a wide code is incomplete.
% HANDOFF item 68 (stage 84 B): causal/stage84/width.txt

% --- add to the bibliography of draft.tex ---------------------------------
% \bibitem[Wang and Isola(2020)]{wang2020alignment}
% T.~Wang and P.~Isola.
% \newblock Understanding contrastive representation learning through alignment
% and uniformity on the hypersphere.
% \newblock \emph{ICML}, 2020.
 after the Background section of draft.tex.
% Needs amsmath, amsthm, booktabs, natbib, xcolor (all loaded by draft.tex).
% The source of every number is in a % comment: items and stages of
% out/burden_feasibility_20261102/HANDOFF_archive_20261127.md (the full record,
% both sessions' copies merged); HANDOFF.md beside it summarises them.
% One new bibliography entry is at the end of this file, commented out.
% =====================================================================
\makeatletter
\@ifundefined{theorem}{\newtheorem{theorem}{Theorem}}{}
\@ifundefined{proposition}{\newtheorem{proposition}{Proposition}}{}
\@ifundefined{lemma}{\newtheorem{lemma}{Lemma}}{}
\makeatother
\providecommand{\R}{\mathbb{R}}
\providecommand{\pending}[1]{{\color{red}\textbf{[PENDING: #1]}}}
\providecommand{\todo}[1]{{\color{blue}\small[TODO: #1]}}

\section{Identifiability at the matched width, and how to reach it}
\label{sec:theory}

Multiview identifiability results promise that a contrastive code of the right
width is a one-to-one function of the content the views share. On our benchmark
that promise is kept in half of the training runs. This section locates the gap
and closes it. We state the result as two steps, each certifying one half of its
conclusion (\S\ref{sec:theory-result}); show that the step certifying
completeness is the one that fails with finite data, and why
(\S\ref{sec:theory-gap}); give a construction that reaches the theorem's
conclusion with a certificate that can be checked on held-out data
(\S\ref{sec:theory-chart}); and state what that certificate does not establish
(\S\ref{sec:theory-check}).

\subsection{Setting}
\label{sec:theory-setting}

Each subject has shared content $c \in C \subset \R^{n_c}$ and, for each view
$v \in \{1, 2\}$ (T1 and FLAIR), a view-specific style $s_v$, drawn
independently given $c$. The views are $x^{(v)} = f_v(c, s_v)$ with $f_v$
smooth and injective. On the fixed-position benchmark the content is seven
anatomical factors (brain volume, ventricle volume, lesion burden, cortical
thickness, temporal volume, left--right asymmetry and sulcal width), so
$n_c = 7$; the style is each view's intensity gain, bias and noise level. Each
view has its own encoder $g_v$, and the two views' codes are compared by the
cosine, that is, on the unit sphere $S^{d-1}$ of a $d$-unit code. With many
negatives, InfoNCE at temperature $\tau$ approaches the alignment
$\mathbb{E}\,\lVert g_1(x^{(1)}) - g_2(x^{(2)}) \rVert^2$ minus the entropy of
each view's code on the sphere \citep{wang2020alignment,zimmermann2021contrastive}.

\subsection{The multiview result is two certificates}
\label{sec:theory-result}

\begin{theorem}[informal; \citealp{vonkugelgen2021self,daunhawer2023identifiability,yao2024multiview}]
\label{thm:block}
Let $g_1$ and $g_2$ minimise the alignment minus the entropy, with the code's
sphere of the content's dimension, $d - 1 = n_c$. Then each view's code is
$g_v(x^{(v)}) = h(c)$ almost surely, for one smooth bijection $h$ of $C$ onto its
image: the code is a function of the content alone, and the content can be
recovered from it.
\end{theorem}
\todo{state the exact theorem and its assumptions as in \citet{yao2024multiview};
here both views observe all of the content, so their partial observability
reduces to the two-view case.}

The proof has two steps, and only the second uses the width.
\begin{enumerate}
\item \emph{Alignment certifies ``only''.} At zero alignment loss
  $g_1(f_1(c, s_1)) = g_2(f_2(c, s_2))$ almost surely. Because $s_1$ and $s_2$
  vary independently given $c$ (with full support, and the maps smooth), neither
  side can depend on its style, so both are one function $h(c)$. This holds at
  any width.
\item \emph{Maximum entropy at the matched width certifies ``all''.} On
  $S^{n_c}$ the entropy is maximal only for the uniform density, which is
  regular. A lemma of \citet{zimmermann2021contrastive} then makes $h$
  one-to-one: a smooth map between simply connected manifolds of the same
  dimension that carries a regular density to a regular density is a
  bijection. \todo{check the lemma's exact statement and number}
\end{enumerate}
Under the cosine the matched width is therefore $d = n_c + 1$ units, eight here;
with an output in $\R^{d}$ rather than on a sphere it would be $n_c$. Above the
matched width, step 2 has nothing to say: an $n_c$-dimensional image in a
higher-dimensional sphere has no finite entropy, so for a wide code the theorem
guarantees only step 1's conclusion.

\subsection{The entropy step fails at the matched width with finite data}
\label{sec:theory-gap}

The entropy in step 2 is never observed. InfoNCE estimates it from batches of
subjects at temperature $\tau$, so it penalises codes that crowd subjects
together at roughly that scale. Consider a code that ignores one direction of
$C$, the variation of one factor's own part, and bends a curved function of the
factors it does hold through the dimension this frees. In the population the
lemma excludes it: its image has dimension $n_c - 1$, a set of measure zero in
$S^{n_c}$, so its density is singular and its entropy is $-\infty$. At the
resolution a batch can see, the bent image can be spread as evenly as a code
that keeps the direction.

Alignment then prefers the curve. A curved function of well-read factors is
computed identically from both views, while the weakest own part (here in the
plane of ventricle size and lesion burden, the variation their causes do not
explain) is determined only partially by each view, so keeping it costs
alignment. The finite-sample optimum can therefore trade the weakest direction
for a curve, and nothing in the objective tells the two apart.

The benchmark shows this trade. % HANDOFF items 46, 56, 58, 61, 63, 64
\begin{itemize}
\item At the matched width (eight units, 2{,}000 training subjects), the code is
  complete\footnote{A code is \emph{complete} when every direction of the seven
  factors' own parts (each factor's variation not explained by its causes) is
  recovered from it with $R^2 \geq .40$ in both views, the ventricle--burden
  plane included. \todo{align with the evaluation section's definition}} in 5
  of 10 training runs that differ only in their seed. % item 61: A100s, seeds 6-15
  Complete codes spend their eighth direction on a curve; incomplete ones spend
  two directions on curves and lose an own-part direction, typically read at
  $R^2$ near zero. % items 46, 56, 58
\item No code carries style at any width: the gain, bias and noise level of
  either view read $R^2 = 0.00$ from codes of 4 to 16 units, linearly and by
  gradient boosting. Step 1 works. % item 63
\item Wider codes put their extra directions into curves. Past the seventh,
  every direction is a joint non-linear function of the seven factors (boosted
  $R^2 \approx .75$) rather than new content, and the codes' intrinsic dimension
  stays at the factors': $6.36$ for 16-unit codes against $6.37$ for the true
  factors on the same subjects (TwoNN). % items 63, 64
\item Nine changes to training at eight units fail to raise the rate: the
  interval at which code batches are rebuilt, the prediction weight, a
  compressed start, a harmonic penalty, the temperature, a projector head at the
  code's width, a wider projector that lets the loss bend the code into sixteen
  outputs (2 of 10 complete), a term penalising what the views still share
  beyond the code (0 of 10), and refitting the head on a frozen backbone.
  % stages 60, 62-64, 70, 71, 71a (item 58); stage 83 (item 69); stage 77 (item 67); stage 80 (item 65)
\item Training longer helps slowly and not always. With 40{,}000 further steps,
  three times the usual length, one of two incomplete codes recovered its lost
  direction (its $R^2$ rose steadily from $.01$ to $.44$), while the other's
  stayed at zero; 10{,}000 further steps had changed neither. The weakest factor
  is learned last.
  % stage 81 (item 70); stage 65 for the 10,000-step continuation
\end{itemize}
The content is identifiable at the matched width, but with finite data the
objective that should find it does not certify completeness.

\subsection{Identify wide, chart narrow}
\label{sec:theory-chart}

The two certificates can be separated. Step 1 does not need the matched width,
and completeness can be certified by something other than entropy. We
therefore (i) train the same objective at a width slightly above the matched
one, where the curves have room and step 1 still holds, and (ii) bring that code
down to the content's dimension with a non-linear chart fitted to rebuild it.

\paragraph{Construction (no labels).}
\begin{enumerate}
\item Train the encoders at a width $d_w > n_c + 1$ (ten units here). By step 1,
  $g_v(x^{(v)}) = \varphi(c)$ for a content-only $\varphi$; that $\varphi$ is
  one-to-one is measured, not guaranteed.
\item In each view, standardise the code, fit a canonical correlation analysis
  between the views on training subjects, keep the directions whose agreement
  on held-out training subjects exceeds $.5$ and a shuffled-pairing null, and
  whiten them. Write $u^{(v)} \in \R^{k}$ for the result.
\item Fit an encoder $\psi : \R^{k} \to \R^{m}$ and a decoder
  $\psi' : \R^{m} \to \R^{k}$ (small MLPs, shared by the views) to minimise
  $\mathbb{E}\,\lVert \psi'(\psi(u^{(v)})) - u^{(v)} \rVert^2$, with early
  stopping on held-out training subjects and the best of five restarts. The
  narrow code is $z^{(v)} = \psi(u^{(v)})$.
\end{enumerate}

\begin{proposition}[chart]
\label{prop:chart}
Let $\varphi : C \to \R^{k}$ be continuous and one-to-one, and let
$\psi'(\psi(u)) = u$ for every $u \in \varphi(C)$. Then $\psi \circ \varphi$ is
continuous and one-to-one. If $m = n_c$, it is a homeomorphism of $C$ onto an
open subset of $\R^{n_c}$, and a diffeomorphism when $\psi$ and $\varphi$ are
smooth with Jacobians of full rank.
\end{proposition}
\begin{proof}
If $\psi(\varphi(c)) = \psi(\varphi(c'))$, applying $\psi'$ gives
$\varphi(c) = \varphi(c')$, so $c = c'$. Continuity is inherited. For
$m = n_c$, a continuous one-to-one map from an open subset of $\R^{n_c}$ into
$\R^{n_c}$ is open by invariance of domain, hence a homeomorphism onto its
image.
\end{proof}

The narrow code is the same function of either view's wide code, so it inherits
step 1's conclusion, and Proposition~\ref{prop:chart} supplies step 2's.
Together they give the conclusion of Theorem~\ref{thm:block} at width $n_c$, one
unit below the cosine's matched width, because the chart's output is not on a
sphere.

Reconstruction succeeds where entropy did not for two reasons. The chart is
fitted to content that is already aligned, so keeping the weakest direction
costs it nothing. And with whitened targets every shared direction counts
equally, while a curve needs no unit of its own: the decoder rebuilds it from
the coordinates of the factors it is a function of. Entropy has neither
property.

\begin{table}[t]
\centering
\caption{Charts of wide codes on the fixed-position benchmark. A code counts when
it is complete. The linear chart (the top principal components of the same
whitened variates) is the best linear compression of the same width.}
\label{tab:chart}
\begin{tabular}{lccc}
\toprule
Wide code $\to$ chart & wide codes & non-linear chart & linear chart\\
\midrule
10 $\to$ 8 units & 4 of 5 & 4 of 4 & 0 of 4\\
10 $\to$ 7 units & 4 of 5 & 4 of 4 & 0 of 4\\
9 $\to$ 8 (code batches every 500 steps) & 3 of 4 & 3 of 3 & 0 of 3\\
9 $\to$ 8 (code batches every 100 steps) & 1 of 4 & 1 of 1 & 0 of 1\\
12 $\to$ 8 & 3 of 3 & 2 of 3 & 0 of 3\\
16 $\to$ 8 & 2 of 3 & 0 of 2 & 0 of 2\\
\midrule
Trained at 8 units (reference) & \multicolumn{3}{c}{5 of 10 complete}\\
\bottomrule
\end{tabular}
\end{table}
% HANDOFF item 66 (stages 82 and 82b); item 61 for the reference row

Wherever its wide code is complete, the chart of a 9- or 10-unit code is
complete: 8 of 8 at eight units and 4 of 4 at seven (Table~\ref{tab:chart}). Its
own-part $R^2$ stays within $.05$ of the wide code's (for example ventricle
$.80 \to .78$ and sulcal width $.74 \to .73$, means over the complete 10-unit
codes and both views). The linear chart is never complete. The charts also keep
their wide codes' causal outputs: average treatment effects estimated with a
smooth learner agree within $.02$, and causal discovery with the code as one
block node recovers its edges to tabular variables with $F_1$ of $.88$--$.94$,
against $.88$--$.96$ for the wide codes and $.98$ for the true own parts.
% item 66, downstream: causal/stage82/tabular_rbf.txt, block_tabular.txt

\subsection{The rebuild check, and what it does not certify}
\label{sec:theory-check}

Proposition~\ref{prop:chart} needs exact rebuilding on $\varphi(C)$. With finite
data, the rebuild error measures how far from it a chart is.

\begin{lemma}[rebuild bound]
\label{lem:rebuild}
For any $\psi$ and $\psi'$ and every direction $j$ of $u$,
$\mathbb{E}\,\mathrm{Var}\bigl(u_j \mid \psi(u)\bigr) \le
\mathbb{E}\,\bigl(u_j - \psi'(\psi(u))_j\bigr)^2$.
\end{lemma}
\begin{proof}
Among all functions of $\psi(u)$, the conditional mean
$\mathbb{E}[u_j \mid \psi(u)]$ has the smallest mean squared error for $u_j$, and
$\psi'(\psi(u))_j$ is one of those functions.
\end{proof}

The share of each whitened shared direction that the narrow code leaves
undetermined is therefore at most its rebuild error, and exact rebuilding gives
the proposition's hypothesis. We estimate the error per direction on test
subjects that neither the canonical analysis, the chart nor its early stopping
saw, with a 95\% bootstrap upper bound over subjects. For the complete 10-unit
codes, the worst direction of the 8-unit charts is at most $.017$ (the mean over
directions at most $.004$), and that of the 7-unit charts at most $.033$.
% item 66, certificate: causal/stage82/cert_u{7,8}{a,b}.txt

Three properties of the check matter.
\begin{itemize}
\item \emph{It is per whitened direction.} An average error, or one in
  unwhitened units, is dominated by high-variance structure, so a faint factor
  can be missing entirely while the average error looks excellent. This is why
  rebuilding the images would be a weak test of whether a faint factor is held.
\item \emph{It is on held-out subjects.} On the fitting subjects near-perfect
  rebuilding is cheap: almost any map is one-to-one on a finite set, and a
  flexible decoder can memorise it.
\item \emph{It certifies information, not readability, and only for continuous
  charts.} A one-to-one code can hold a factor in a form a probe cannot read,
  and without continuity a single real number can encode several factors, so
  the dimension count in Proposition~\ref{prop:chart} rests on continuity. The
  own-part readouts check readability directly.
\end{itemize}

The check does not certify the wide code. It compares the chart with its wide
code, not the wide code with the content, and a chart cannot add what its wide
code lacks: the one incomplete 10-unit code is rebuilt to within $.036$, and its
chart is still incomplete. The construction is therefore as reliable as the wide
code's completeness, 4 of 5 at ten units. Wider codes chart worse with finite
data (2 of 3 at twelve units, 0 of 2 at sixteen): every extra curve direction
counts fully in the whitened target, so the bottleneck again chooses between a
curve and the weakest factor. A slightly wide code is the right parent. Two
label-free steps complete the construction: choosing a complete code among
several wide ones, and setting the chart's width when $n_c$ is unknown, as on
ADNI.

\paragraph{Choosing a complete wide code without labels.} Every content-only
code is a function of the content, so a complete code determines every other
run's code, while a code that misses a part cannot determine one that holds it.
We therefore train several seeds and score each code by how much of the others'
shared content it predicts. For a candidate $A$ and a reference $B$, each view
of $A$ predicts the other view's shared variates of $B$ (so view-specific
variation cannot count), against $B$'s own view as the ceiling; $q(A \mid B)$ is
the smallest share, over the directions of $B$'s whitened variates, of what $B$'s
own view explains that $A$ also explains. A code's score is its minimum $q$ over
the other seeds, and we keep the code with the highest score. At nine and ten
units the complete code scores higher in 10 of 10 complete--incomplete pairs,
where the backbone-leftover criterion of \S\todo{ref} gets 7, and among any
three seeds the top score is complete in every triple that contains a complete
code. At eight units 68 of 69 pairs are right. Because the score reads a code
against other runs' codes rather than against its own backbone, a backbone that
carries less cannot lower it: the 20 codes of a training variant that degraded
its backbone, all incomplete, score below every complete reference code.
% HANDOFF item 68 (stage 84 A); the variant is stage 77 (item 67); rho1 at 9-10 units: item 60

\paragraph{Setting the chart's width without labels.} The rebuild check also
gives the width. Charting each complete 10-unit code at decreasing widths, the
worst-direction error (95\% bound) stays at or below $.033$ down to seven units
and jumps to $.17$--$.31$ at six, where every chart loses an own part (sulcal
width first). The smallest width at which the chart rebuilds the wide code to
within $.10$ is therefore $n_c = 7$ for every complete code, with no labels and
no calibration borrowed from known latents. The one incomplete 10-unit code
shows no jump at seven and is rebuilt at six units to within $.08$: its shared
part is one dimension short, as its own parts confirm, so a width one below the
others is a second sign that a wide code is incomplete.
% HANDOFF item 68 (stage 84 B): causal/stage84/width.txt

% --- add to the bibliography of draft.tex ---------------------------------
% \bibitem[Wang and Isola(2020)]{wang2020alignment}
% T.~Wang and P.~Isola.
% \newblock Understanding contrastive representation learning through alignment
% and uniformity on the hypersphere.
% \newblock \emph{ICML}, 2020.
